\clearpage\thispagestyle{empty}
\pdfbookmark[0]{Quellen zu Borda’s Kreis}{bordacircle1818-sources}
\begin{center}\textbf{Quellen zu Borda’s Kreis}\end{center}
\noindent Diese KI-Übertragung ins heutige Deutsch folgt Carl Friedrich Gauß’ \textit{Borda’s Kreis} in den \href{https://gdz.sub.uni-goettingen.de/id/PPN236007467}{\textit{Werken}, Band~VI} (Göttingen, 1874), S.~404--409. Der Erstabdruck erschien als Beitrag XXIV, \textit{Ueber einige Berichtigungen an Bordaischen Wiederholungskreisen}, in der von Bernhard von Lindenau und Johann Gottlieb Friedrich Bohnenberger herausgegebenen \textit{Zeitschrift für Astronomie und verwandte Wissenschaften}, Band~V, März und April 1818, S.~\href{https://api.digitale-sammlungen.de/iiif/image/v2/bsb10061007_00208/full/1600,/0/default.jpg}{198}--\href{https://api.digitale-sammlungen.de/iiif/image/v2/bsb10061007_00221/full/1600,/0/default.jpg}{211}. Die Links zum Erstabdruck führen zum Exemplar der Bayerischen Staatsbibliothek. In der Quellenzeile der Werkausgabe steht \textit{Liadenau} für Lindenau.
\par\smallskip\noindent Auf S.~\href{https://api.digitale-sammlungen.de/iiif/image/v2/bsb10061007_00210/full/1600,/0/default.jpg}{200} definiert der Erstabdruck $M$ und $N$ mit der zweiten Achse als gemeinsamer Achse. In der Werkausgabe ist es die dritte. Auf S.~\href{https://api.digitale-sammlungen.de/iiif/image/v2/bsb10061007_00214/full/1600,/0/default.jpg}{204} gibt der Erstabdruck $\cos I=\sin a\sin b$ und einen Ausdruck für den Tangens von $\eta$ mit $\sin a\sin b$ an. In der Werkausgabe steht in der Kosinusformel das Produkt der Tangenswerte von $b$ und $c$ mit einem Minuszeichen. Für den Tangens stehen dort zwei Wurzelausdrücke. Auf S.~\href{https://api.digitale-sammlungen.de/iiif/image/v2/bsb10061007_00217/full/1600,/0/default.jpg}{207} nennt der Erstabdruck bei der horizontalen Einstellung $a$ und $b$. Die Werkausgabe nennt $c$ und $b$. Die Übertragung folgt der Werkausgabe und übernimmt ihre Wurzelzeichen ohne Überstrich und die verschiedenen Formen für $\upsilon$ und $V$. KI-Übertragung und KI-Satz.
\clearpage
